\chapter{HITL Audit Log}
\label{chap:HITL_Audit_Log}

\section{Purpose}

The HITL Audit Log records all significant human interventions in the CME:
\begin{itemize}
    \item approvals and rejections,
    \item SIGIL invocations,
    \item policy changes,
    \item manual overrides,
    \item EAC and CAR decisions.
\end{itemize}

It is the primary accountability and forensics tool for governance.

\section{Entry Schema}

An audit entry $A$ is:
\[
A = \langle id,\, time,\, actor,\, action,\, target,\, rationale,\, outcome \rangle
\]

Where:
\begin{itemize}
    \item id: unique audit entry ID,
    \item time: timestamp,
    \item actor: human or institutional identity,
    \item action: what was done (approve / reject / override / etc.),
    \item target: Engram, policy, SIGIL, or system context,
    \item rationale: short natural-language reason,
    \item outcome: success / failure / partial / escalated.
\end{itemize}

\section{Integrity Requirements}

\begin{itemize}
    \item entries must be append-only,
    \item no deletion or in-place modification,
    \item corrections create new entries referencing prior ones,
    \item logs must be tamper-evident (cryptographic integrity checks).
\end{itemize}

\section{Access Control}

\begin{itemize}
    \item operators may view their own entries and aggregated summaries,
    \item governance roles may view wider slices,
    \item EAC may have full read access in sensitive cases,
    \item write access is only through harness logging, never direct edits.
\end{itemize}

\section{Usage}

The audit log supports:
\begin{itemize}
    \item investigations after incidents,
    \item pattern analysis of high-risk actions,
    \item evaluation of governance processes,
    \item trust-building with stakeholders through traceability.
\end{itemize}

\section{Summary}

The HITL Audit Log is the memory of human governance. Without it, there is
no way to reliably reconstruct who did what, when, and why in the life of
the CME.
